% =============================================================================
% SPÖ Inflation Strategy — FehrAdvice Presentation Deck
% =============================================================================
% Format: 16:9 Beamer (FehrAdvice Theme)
% Version: 1.1 | 2026-02-06
% Kompilierung: pdflatex DECK_inflation_ordnung_2026-02-06.tex
% Theme SSOT: templates/beamerthemefehradvice.sty
% =============================================================================

\documentclass[aspectratio=169, 11pt, xcolor={table}]{beamer}

% FA Theme laden
\usepackage{beamerthemefehradvice}

\usepackage[ngerman]{babel}

% --- Metadata ---
\title{Wirtschaftliche Ordnung f\"ur \"Osterreich}
\subtitle{Inflationsstrategie -- Evidenzbasierte Kommunikation}
\author{Dr.\ Gerhard Fehr}
\institute{FehrAdvice \& Partners AG}
\date{Februar 2026}

\begin{document}

% =============================================================================
% SLIDE 0: DECKBLATT
% =============================================================================

\begin{frame}[plain]
    \titlepage
\end{frame}

% =============================================================================
% SLIDE 1: DAS PROBLEM
% =============================================================================

\begin{frame}{Das Problem: Inflation als Vertrauenskrise}
\framesubtitle{Warum Zahlen allein nicht reichen}

\begin{itemize}
    \item Die Inflation in \"Osterreich ist von 11\,\% auf 2\,\% gesunken --
          doch \textbf{73\,\% der Bev\"olkerung} sp\"uren keine Entlastung.
    \item Der \textbf{Wahrnehmungs-Gap} zwischen Statistik und Erleben
          ist das zentrale Problem f\"ur die Regierungskommunikation.
    \item Populistische Narrative (\guillemotleft Alles wird teurer\guillemotright)
          dominieren, weil sie \textbf{System~1} (intuitiv, emotional) ansprechen.
\end{itemize}

\vskip8pt
\takeaway{Wir m\"ussen die Wahrnehmung gestalten, nicht nur die Zahlen kommunizieren.}

\end{frame}

% =============================================================================
% SLIDE 2: DIE EVIDENZ
% =============================================================================

\begin{frame}[shrink=2]{Die Evidenz: Was die Daten zeigen}
\framesubtitle{Drei Schl\"usselkennzahlen zur Preisentwicklung}

\begin{columns}[T]
    \begin{column}{0.32\textwidth}
        \bigstat{11\,\% $\rightarrow$ 2\,\%}{Inflationsrate\\(Peak $\rightarrow$ aktuell)}{fadarkblue}
    \end{column}
    \begin{column}{0.32\textwidth}
        \bigstat{--0,7\,\%}{Preisr\"uckgang\\J\"anner 2026}{faorange}
    \end{column}
    \begin{column}{0.32\textwidth}
        \bigstat{--4,9\,\%}{Energiepreise\\($\approx$ 300\,\texteuro{}/Jahr)}{fasuccessgreen}
    \end{column}
\end{columns}

\vskip4pt
\insightbox{Behavioral Insight}{%
    \small Der monatliche Preisr\"uckgang von 0,7\,\% im J\"anner 2026 ist der
    \textbf{st\"arkste seit J\"anner 2019}. Reframing
    (von \guillemotleft Preise steigen langsamer\guillemotright{} zu
    \guillemotleft Preise sinken\guillemotright) ist kommunikativ entscheidend.
}

\end{frame}

% =============================================================================
% SLIDE 3: DER RAHMEN (Ordnung vs. Chaos)
% =============================================================================

\begin{frame}[shrink=5]{Der Rahmen: Ordnung statt Spalten}
\framesubtitle{Frame-Wettbewerb als prim\"are Intervention (52\,\% Impact)}

\begin{columns}[T]
    \begin{column}{0.46\textwidth}
        \begin{tcolorbox}[
            colback=falightgray, colframe=famedgray,
            title={\color{fadarkgray}\bfseries\small Bisheriges Narrativ (FP\"O)},
            colbacktitle=falinegray,
            left=4pt, right=4pt, top=2pt, bottom=2pt
        ]
            \small
            \begin{itemize}\setlength{\itemsep}{0pt}
                \item \guillemotleft Alles wird teurer\guillemotright
                \item \guillemotleft Die Regierung tut nichts\guillemotright
                \item \guillemotleft System ist kaputt\guillemotright
            \end{itemize}
            \vskip2pt
            {\small\color{fadangerred}\bfseries $\rightarrow$ Chaos-Frame}
        \end{tcolorbox}
    \end{column}
    \begin{column}{0.46\textwidth}
        \begin{tcolorbox}[
            colback=fadarkblue!5!white, colframe=fadarkblue,
            title={\color{fawhite}\bfseries\small Neues Narrativ (SP\"O)},
            colbacktitle=fadarkblue,
            left=4pt, right=4pt, top=2pt, bottom=2pt
        ]
            \small
            \begin{itemize}\setlength{\itemsep}{0pt}
                \item \guillemotleft Preise sinken messbar\guillemotright
                \item \guillemotleft Schritt f\"ur Schritt -- mit Ordnung\guillemotright
                \item \guillemotleft Klare Regeln statt leere Versprechen\guillemotright
            \end{itemize}
            \vskip2pt
            {\small\color{fasuccessgreen}\bfseries $\rightarrow$ Ordnung-Frame}
        \end{tcolorbox}
    \end{column}
\end{columns}

\vskip4pt
\takeaway{Frame-Wettbewerb dominiert mit 52\,\% Impact -- IMMER vor Zahlen einsetzen.}

\end{frame}

% =============================================================================
% SLIDE 4: DER VERGLEICH (Österreich vs. Ungarn)
% =============================================================================

\begin{frame}[shrink=3]{Der Vergleich: Ordnung wirkt}
\framesubtitle{\"Osterreich vs.\ Ungarn -- Aktive Markt-Ordnung vs.\ Populismus}

\vskip4pt

\compbar{\"Osterreich (SP\"O)}{2,0\,\%}{0.25}{fasuccessgreen}

\vskip8pt

\compbar{Ungarn (Orb\'an)}{4,0\,\%}{0.50}{fadangerred}

\vskip10pt

\begin{columns}[T]
    \begin{column}{0.46\textwidth}
        \begin{center}
            {\Large\bfseries\color{fasuccessgreen}$\checkmark$}\\[2pt]
            {\small\bfseries Aktive Markt-Ordnung}\\
            {\footnotesize\color{famedgray} MwSt-Senkung, Mietpreisbremse,\\Anti-Shrinkflation-Gesetz}
        \end{center}
    \end{column}
    \begin{column}{0.46\textwidth}
        \begin{center}
            {\Large\bfseries\color{fadangerred}$\times$}\\[2pt]
            {\small\bfseries Populistische Preispolitik}\\
            {\footnotesize\color{famedgray} Preisdeckel ohne Marktmechanismus,\\\"okonomische Instabilit\"at}
        \end{center}
    \end{column}
\end{columns}

\vskip4pt
\takeaway{Vergleich mit Ungarn macht Ordnungs-Vorteil konkret und messbar.}

\end{frame}

% =============================================================================
% SLIDE 5: DIE MASSNAHMEN (3 Schritte)
% =============================================================================

\begin{frame}[shrink=5]{Die Massnahmen: Drei Schritte zur Entlastung}
\framesubtitle{Konkrete Policy-Interventionen mit messbarer Wirkung}

\vskip2pt

\numpoint{1}{1,4 Mrd.\ \texteuro{} MwSt-Senkung}\\[1pt]
{\small\color{famedgray}\hspace{1cm} Auf Lebensmittel des t\"aglichen Bedarfs -- sp\"urbar beim n\"achsten Einkauf.}

\vskip8pt

\numpoint{2}{Mietpreisbremse f\"ur 300.000 Haushalte}\\[1pt]
{\small\color{famedgray}\hspace{1cm} Indexierung gedeckelt -- Schutz vor unkontrollierten Mieterh\"ohungen.}

\vskip8pt

\numpoint{3}{Anti-Shrinkflation-Gesetz}\\[1pt]
{\small\color{famedgray}\hspace{1cm} Transparenzpflicht bei Mengenreduktion -- faire Preise statt versteckter Teuerung.}

\vskip6pt

\insightbox{Ehrlichkeits-Linie (BehavEcon)}{%
    \small\guillemotleft Gegen die Teuerung gewinnt man nie ganz.
    Aber Schritt f\"ur Schritt -- mit Ordnung.\guillemotright\\[1pt]
    {\small\color{famedgray} Senkt Erwartungen realistisch $\rightarrow$ \"Uberperformance wird positiv wahrgenommen.}
}

\end{frame}

% =============================================================================
% SLIDE 6: DEPLOYMENT & EMPFEHLUNG
% =============================================================================

\begin{frame}[shrink=5]{Empfehlung: Deployment nach Kontext}
\framesubtitle{Interventions-Mix nach Anlass und Zielgruppe}

\vskip1pt
{\footnotesize
\begin{tabular}{p{2.5cm}p{4.2cm}p{4.2cm}}
    \rowcolor{fadarkblue}
    {\color{fawhite}\bfseries Anlass} &
    {\color{fawhite}\bfseries Lead-Intervention} &
    {\color{fawhite}\bfseries Infografik-Format} \\
    \rowcolor{falightgray}
    TV-Interview & Zahlen-Push + Ordnung-Frame & -- \\
    Social Media & Killer-Zahlen + Infografik & Square (IG/FB) \\
    \rowcolor{falightgray}
    Parteitag & Werte + Abgrenzung & Landscape (Screen) \\
    Wahlkampf Strasse & Pers\"onliche Geschichten & Portrait (Flyer) \\
    \rowcolor{falightgray}
    Pressekonferenz & Zahlen + Vergleich AT/HU & Landscape (Handout) \\
\end{tabular}
}

\vskip4pt

\begin{columns}[T]
    \begin{column}{0.46\textwidth}
        \insightbox{Wechselw\"ahler Mitte}{%
            \small\textbf{Do:} Fakten, Zahlen, Ordnung-Frame\\
            \textbf{Don't:} Parteipolitik, FP\"O-Angriffe
        }
    \end{column}
    \begin{column}{0.46\textwidth}
        \insightbox{\"OVP-Entt\"auschte}{%
            \small\textbf{Do:} Ordnung, Kompetenz, Pragmatismus\\
            \textbf{Don't:} \"OVP-Kritik, Links-Rhetorik
        }
    \end{column}
\end{columns}

\end{frame}

% =============================================================================
% SLIDE 7: SCHLUSSSEITE
% =============================================================================

\closingslide
    {Vielen Dank f\"ur \mbox{Ihre} Aufmerksamkeit!}
    {Dr.\ Gerhard Fehr}
    {fehr@fehradvice.com}

\end{document}
